% ===========================================================================
% DANH MỤC TÀI LIỆU THAM KHẢO
% ===========================================================================
\chapter*{DANH MỤC TÀI LIỆU THAM KHẢO}
\addcontentsline{toc}{chapter}{DANH MỤC TÀI LIỆU THAM KHẢO}
\markboth{DANH MỤC TÀI LIỆU THAM KHẢO}{DANH MỤC TÀI LIỆU THAM KHẢO}

% Danh mục tài liệu chứa nhiều tên riêng và tiêu đề tiếng Anh mà bộ quy tắc ngắt
% âm tiết tiếng Việt không xử lý được. Nới lỏng tiêu chí xuống dòng trong phạm vi
% chương này để tránh tràn lề; chấp nhận khoảng trắng giữa từ thưa hơn đôi chút.
% Danh mục tài liệu được căn lề trái (không căn đều hai bên).
%
% Lý do: định dạng XeLaTeX đang dùng không nạp sẵn bộ quy tắc ngắt âm tiết
% (kiểm chứng bằng \showhyphens: không có điểm ngắt nào được đề xuất). Với phần
% chữ tiếng Việt điều này không gây vấn đề vì tiếng Việt vốn không ngắt âm tiết;
% nhưng các tiêu đề và tên riêng tiếng Anh dài trong danh mục tài liệu sẽ không
% ngắt được, buộc TeX phải giãn khoảng trắng rất thưa để căn đều hai bên.
%
% Căn lề trái là cách trình bày chuẩn và phổ biến cho danh mục tài liệu tham
% khảo, đồng thời loại bỏ hoàn toàn hiện tượng dòng thưa nói trên.
\begingroup
\raggedright
\emergencystretch=2.5em

\section*{Tài liệu tham khảo tiếng Việt}

\begin{enumerate}[label={[\arabic*]}, ref=\arabic*, leftmargin=*, labelsep=0.6em,
                  itemsep=5pt]

\item\label{ref:tt15} Ngân hàng Nhà nước Việt Nam (2023), \textit{Thông tư số
15/2023/TT-NHNN quy định về hoạt động thông tin tín dụng của Ngân hàng Nhà nước
Việt Nam}, Hà Nội.

\item\label{ref:tt11} Ngân hàng Nhà nước Việt Nam (2021), \textit{Thông tư số
11/2021/TT-NHNN quy định về phân loại tài sản có, mức trích, phương pháp trích
lập dự phòng rủi ro và việc sử dụng dự phòng để xử lý rủi ro trong hoạt động của
tổ chức tín dụng, chi nhánh ngân hàng nước ngoài}, Hà Nội.

\item\label{ref:nhnn-ml} Ngân hàng Nhà nước Việt Nam (2023), \textit{Machine
Learning trong quản trị rủi ro tín dụng của ngân hàng}, Tạp chí Ngân hàng.

\item\label{ref:thu2025} Huỳnh Thị Ái Thu (2025), \textit{Ứng dụng học máy trong
xây dựng mô hình cảnh báo sớm rủi ro tín dụng khách hàng cá nhân tại ngân hàng
thương mại Việt Nam}, Luận văn Thạc sĩ, Trường Đại học Kinh tế Thành phố Hồ Chí
Minh.

\item\label{ref:thao2025} Nguyễn Thị Phương Thảo (2025), \textit{Xây dựng mô hình
chấm điểm xếp hạng tín dụng khách hàng doanh nghiệp tại các ngân hàng thương mại
ứng dụng các mô hình học máy hiện đại}, Luận văn Thạc sĩ, Trường Đại học Kinh tế
Thành phố Hồ Chí Minh.

\item\label{ref:sang2021} Nguyễn Minh Sáng, Nguyễn Thị Hồng Vinh (2021),
``Lược khảo các nghiên cứu về ứng dụng công nghệ phân tích dữ liệu lớn và trí tuệ
nhân tạo trong quản trị rủi ro tín dụng của ngân hàng'', \textit{Tạp chí Kinh tế
và Ngân hàng châu Á}, số 186.

\item\label{ref:hoa2025} Nguyễn Thị Hoa, Nguyễn Thị Thu (2025), ``Ứng dụng phân
tích dữ liệu lớn và trí tuệ nhân tạo trong tối ưu hóa quản trị rủi ro tín dụng
tại các ngân hàng số ở Việt Nam'', \textit{Tạp chí Khoa học Trường Đại học Mở Hà
Nội}, số 2A.

\end{enumerate}

\section*{Tài liệu tham khảo tiếng Anh}

\begin{enumerate}[label={[\arabic*]}, ref=\arabic*, leftmargin=*, labelsep=0.6em,
                  itemsep=5pt, start=8]

\item\label{ref:lessmann2015} Lessmann, S., Baesens, B., Seow, H.-V., \& Thomas,
L. C. (2015), ``Benchmarking state-of-the-art classification algorithms for credit
scoring: An update of research'', \textit{European Journal of Operational
Research}, 247(1), pp. 124--136.

\item\label{ref:shi2022} Shi, S., Tse, R., Luo, W., D'Addona, S., \& Pau, G.
(2022), ``Machine learning-driven credit risk: A systemic review'', \textit{Neural
Computing and Applications}, 34, pp. 14327--14339.

\item\label{ref:onay2018} Onay, C., \& Öztürk, E. (2018), ``A review of credit
scoring research in the age of Big Data'', \textit{Journal of Financial Regulation
and Compliance}, 26(3), pp. 382--405.

\item\label{ref:ayari2026} Ayari, H., Guetari, R., \& Kraïem, N. (2026),
``Machine learning powered financial credit scoring: A systematic literature
review'', \textit{Artificial Intelligence Review}, 59, 13.

\item\label{ref:data2023} ``Machine Learning for Credit Risk Prediction: A
Systematic Literature Review'' (2023), \textit{Data}, 8(11), 169.

\item\label{ref:cortes1995} Cortes, C., \& Vapnik, V. (1995), ``Support-Vector
Networks'', \textit{Machine Learning}, 20(3), pp. 273--297.

\item\label{ref:vapnik1998} Vapnik, V. N. (1998), \textit{Statistical Learning
Theory}, John Wiley \& Sons.

\item\label{ref:williams2001} Williams, C. K. I., \& Seeger, M. (2001), ``Using
the Nyström Method to Speed Up Kernel Machines'', \textit{Advances in Neural
Information Processing Systems}, 13, pp. 682--688.

\item\label{ref:platt1999} Platt, J. C. (1999), ``Probabilistic Outputs for
Support Vector Machines and Comparisons to Regularized Likelihood Methods'',
\textit{Advances in Large Margin Classifiers}, MIT Press, pp. 61--74.

\item\label{ref:breiman1984} Breiman, L., Friedman, J. H., Olshen, R. A., \&
Stone, C. J. (1984), \textit{Classification and Regression Trees}, Wadsworth.

\item\label{ref:hastie2009} Hastie, T., Tibshirani, R., \& Friedman, J. (2009),
\textit{The Elements of Statistical Learning: Data Mining, Inference, and
Prediction}, 2nd edition, Springer.

\item\label{ref:mitchell1997} Mitchell, T. M. (1997), \textit{Machine Learning},
McGraw-Hill.

\item\label{ref:chen2016} Chen, T., \& Guestrin, C. (2016), ``XGBoost: A Scalable
Tree Boosting System'', \textit{Proceedings of the 22nd ACM SIGKDD International
Conference on Knowledge Discovery and Data Mining}, pp. 785--794.

\item\label{ref:ke2017} Ke, G., Meng, Q., Finley, T., Wang, T., Chen, W., Ma, W.,
Ye, Q., \& Liu, T.-Y. (2017), ``LightGBM: A Highly Efficient Gradient Boosting
Decision Tree'', \textit{Advances in Neural Information Processing Systems}, 30,
pp. 3146--3154.

\item\label{ref:prokhorenkova2018} Prokhorenkova, L., Gusev, G., Vorobev, A.,
Dorogush, A. V., \& Gulin, A. (2018), ``CatBoost: unbiased boosting with
categorical features'', \textit{Advances in Neural Information Processing
Systems}, 31, pp. 6638--6648.

\item\label{ref:breiman2001} Breiman, L. (2001), ``Random Forests'',
\textit{Machine Learning}, 45(1), pp. 5--32.

\item\label{ref:friedman2001} Friedman, J. H. (2001), ``Greedy function
approximation: A gradient boosting machine'', \textit{The Annals of Statistics},
29(5), pp. 1189--1232.

\item\label{ref:lundberg2017} Lundberg, S. M., \& Lee, S.-I. (2017), ``A Unified
Approach to Interpreting Model Predictions'', \textit{Advances in Neural
Information Processing Systems}, 30, pp. 4765--4774.

\item\label{ref:lundberg2020} Lundberg, S. M., Erion, G., Chen, H., DeGrave, A.,
Prutkin, J. M., Nair, B., Katz, R., Himmelfarb, J., Bansal, N., \& Lee, S.-I.
(2020), ``From local explanations to global understanding with explainable AI for
trees'', \textit{Nature Machine Intelligence}, 2(1), pp. 56--67.

\item\label{ref:zadrozny2002} Zadrozny, B., \& Elkan, C. (2002), ``Transforming
classifier scores into accurate multiclass probability estimates'',
\textit{Proceedings of the 8th ACM SIGKDD International Conference on Knowledge
Discovery and Data Mining}, pp. 694--699.

\item\label{ref:niculescu2005} Niculescu-Mizil, A., \& Caruana, R. (2005),
``Predicting good probabilities with supervised learning'', \textit{Proceedings of
the 22nd International Conference on Machine Learning}, pp. 625--632.

\item\label{ref:bergstra2012} Bergstra, J., \& Bengio, Y. (2012), ``Random Search
for Hyper-Parameter Optimization'', \textit{Journal of Machine Learning
Research}, 13, pp. 281--305.

\item\label{ref:siddiqi2017} Siddiqi, N. (2017), \textit{Intelligent Credit
Scoring: Building and Implementing Better Credit Risk Scorecards}, 2nd edition,
John Wiley \& Sons.

\item\label{ref:thomas2017} Thomas, L. C., Crook, J., \& Edelman, D. (2017),
\textit{Credit Scoring and Its Applications}, 2nd edition, Society for Industrial
and Applied Mathematics.

\item\label{ref:baesens2016} Baesens, B., Roesch, D., \& Scheule, H. (2016),
\textit{Credit Risk Analytics: Measurement Techniques, Applications, and Examples
in SAS}, John Wiley \& Sons.

\item\label{ref:bcbs2006} Basel Committee on Banking Supervision (2006),
\textit{International Convergence of Capital Measurement and Capital Standards: A
Revised Framework}, Bank for International Settlements.

\item\label{ref:sculley2015} Sculley, D., Holt, G., Golovin, D., Davydov, E.,
Phillips, T., Ebner, D., Chaudhary, V., Young, M., Crespo, J.-F., \& Dennison, D.
(2015), ``Hidden Technical Debt in Machine Learning Systems'', \textit{Advances in
Neural Information Processing Systems}, 28, pp. 2503--2511.

\item\label{ref:gama2014} Gama, J., Žliobaitė, I., Bifet, A., Pechenizkiy, M., \&
Bouchachia, A. (2014), ``A survey on concept drift adaptation'', \textit{ACM
Computing Surveys}, 46(4), pp. 1--37.

\item\label{ref:chawla2002} Chawla, N. V., Bowyer, K. W., Hall, L. O., \&
Kegelmeyer, W. P. (2002), ``SMOTE: Synthetic Minority Over-sampling Technique'',
\textit{Journal of Artificial Intelligence Research}, 16, pp. 321--357.

\item\label{ref:elkan2001} Elkan, C. (2001), ``The Foundations of Cost-Sensitive
Learning'', \textit{Proceedings of the 17th International Joint Conference on
Artificial Intelligence}, pp. 973--978.

\end{enumerate}
\endgroup
